%! AP Statistics — Unit 1, Lesson 1.9 — Warm-Up
\documentclass[10pt]{article}
\usepackage{apstats-article}
\usepackage{apstats-boxes}

%=============================================================================
\begin{document}

\pageheader{Unit 1, Lesson 1.9}{Warm-Up}
\namedateperiod

\vspace{0.16in}

\noindent In Lessons 1.7--1.8 you built and read boxplots for a single group.
Today you will compare two groups on the \textit{same} display.
The parallel boxplots below show \textbf{ACT composite scores} for two groups of
students: those who did \textit{not} take a test-prep course (\textbf{No Prep})
and those who did (\textbf{Test Prep}).

\vspace{0.18in}

\begin{center}
\begin{tikzpicture}[scale=1.0]
  % Axis: 14 to 36, x = (v - 14) * 0.30, label every 2
  \draw[thick] (0.0, 0) -- (6.80, 0);
  \foreach \v/\x in {14/0.00, 16/0.60, 18/1.20, 20/1.80, 22/2.40,
                      24/3.00, 26/3.60, 28/4.20, 30/4.80, 32/5.40, 34/6.00, 36/6.60} {
    \draw (\x, -0.12) -- (\x, 0.12);
    \node[below, font=\small] at (\x, -0.12) {\v};
  }
  \node[below, font=\small\itshape] at (3.30, -0.56) {ACT composite score};

  % No Prep: min=15(0.30), Q1=19(1.50), M=22(2.40), Q3=25(3.30), max=29(4.50)
  % y_center = 1.10, box from 0.85 to 1.35
  \node[left, font=\small] at (0.0, 1.10) {No Prep};
  \draw[thick, navy] (0.30, 1.10) -- (1.50, 1.10);          % left whisker
  \draw[thick, navy] (1.50, 0.85) rectangle (3.30, 1.35);   % box
  \draw[thick, navy] (2.40, 0.85) -- (2.40, 1.35);          % median
  \draw[thick, navy] (3.30, 1.10) -- (4.50, 1.10);          % right whisker
  \draw[thick, navy] (0.30, 0.95) -- (0.30, 1.25);          % left cap
  \draw[thick, navy] (4.50, 0.95) -- (4.50, 1.25);          % right cap

  % Test Prep: min=20(1.80), Q1=24(3.00), M=27(3.90), Q3=30(4.80), max=35(6.30)
  % y_center = 0.35, box from 0.10 to 0.60
  \node[left, font=\small] at (0.0, 0.35) {Test Prep};
  \draw[thick, navy] (1.80, 0.35) -- (3.00, 0.35);          % left whisker
  \draw[thick, navy] (3.00, 0.10) rectangle (4.80, 0.60);   % box
  \draw[thick, navy] (3.90, 0.10) -- (3.90, 0.60);          % median
  \draw[thick, navy] (4.80, 0.35) -- (6.30, 0.35);          % right whisker
  \draw[thick, navy] (1.80, 0.20) -- (1.80, 0.50);          % left cap
  \draw[thick, navy] (6.30, 0.20) -- (6.30, 0.50);          % right cap
\end{tikzpicture}
\end{center}

\vspace{0.22in}

\begin{enumerate}[leftmargin=1.5em, itemsep=10pt]

  \item Identify the \textbf{median} for each group from the boxplot.\\[6pt]
        No Prep median: \blank{2cm} \quad Test Prep median: \blank{2cm}\\[4pt]
        Which group scored higher on average? \blank{5cm}

  \item Find the \textbf{IQR} for each group.\\[6pt]
        No Prep: $Q_1 = \blank{1.8cm}$, $Q_3 = \blank{1.8cm}$,
        IQR $= \blank{2cm}$ points\\[4pt]
        Test Prep: $Q_1 = \blank{1.8cm}$, $Q_3 = \blank{1.8cm}$,
        IQR $= \blank{2cm}$ points\\[4pt]
        Which group's scores are more spread out? \blank{5cm}

  \item Do the two distributions \textit{overlap}? What range of scores do
        both groups share?\\[6pt]
        \writeline

  \item Write one complete \textbf{AP comparative statement} comparing the
        centers. Use the sentence frame: ``The median ACT score for
        [Group A] was [value], which is [higher/lower] than\ldots''\\[6pt]
        \writeline\\[1pt]\writeline

\end{enumerate}

\vspace{0.14in}

\begin{tcolorbox}[colback=greenbg, colframe=greenacc, boxrule=0.8pt, arc=2mm,
    left=4mm, right=4mm, top=3mm, bottom=3mm]
\small
\begin{itemize}[leftmargin=1.3em, itemsep=5pt]
  \item Today you will learn how to use the SOCS framework to compare
        \textit{two groups} on the same display --- not just describe one.
  \item You will also practice a back-to-back stemplot, a second display
        for comparing small datasets.
  \item Key focus: AP scoring requires \textbf{explicit} comparative
        language. Describing each group separately does not earn credit.
\end{itemize}
\end{tcolorbox}

\vspace{0.12in}

\noindent\textcolor{slate}{\small One question you have going into today's lesson,
or something you're not sure about yet:}\\[8pt]
\writeline

\end{document}
